pdfoutput=1
\pdfoutput=1
\documentclass[12pt,a4paper]{article}

% ============================================================
%  Q-Theorem — Quantum SCX Theorem:
%  Quantum No-Cloning Strengthening of Theorem 3,
%  Entanglement Audit Advantage, and the
%  Quantum Cercis Uncertainty Relation
%  arXiv-ready · English · Standalone
% ============================================================

\usepackage[T1]{fontenc}
\usepackage{amsmath,amssymb,amsfonts}
\usepackage{mathtools}
\usepackage{amsthm}
\usepackage{bm}
\usepackage{graphicx}
\usepackage{xcolor}
\usepackage{hyperref}
\usepackage{booktabs}
\usepackage{enumitem}
\usepackage[numbers]{natbib}
\usepackage{braket}

% ---------- Theorem environments ----------
\newtheorem{theorem}{Theorem}[section]
\newtheorem{lemma}[theorem]{Lemma}
\newtheorem{proposition}[theorem]{Proposition}
\newtheorem{corollary}[theorem]{Corollary}
\newtheorem{definition}[theorem]{Definition}
\newtheorem{remark}[theorem]{Remark}
\newtheorem{assumption}[theorem]{Assumption}
\newtheorem{conjecture}[theorem]{Conjecture}
\newtheorem{openproblem}[theorem]{Open Problem}

% ---------- Honesty critique environment ----------
\newenvironment{critique}[1][Honest Critique]
  {\medskip\noindent\textbf{#1.}\quad\itshape\small}
  {\medskip}

% ---------- Status tags ----------
\newcommand{\rigorous}{\textcolor{green!50!black}{\small\textsf{[Rigorous]}}}
\newcommand{\conditionallyrigorous}{\textcolor{orange}{\small\textsf{[Conditionally Rigorous]}}}
\newcommand{\openquest}{\textcolor{red!70!black}{\small\textsf{[Open]}}}

% ---------- Macros ----------
\newcommand{\R}{\mathbb{R}}
\newcommand{\C}{\mathbb{C}}
\newcommand{\E}{\mathbb{E}}
\newcommand{\V}{\mathbb{V}}
\newcommand{\I}{\mathbb{I}}
\newcommand{\cX}{\mathcal{X}}
\newcommand{\cY}{\mathcal{Y}}
\newcommand{\cD}{\mathcal{D}}
\newcommand{\cH}{\mathcal{H}}
\newcommand{\cL}{\mathcal{L}}
\newcommand{\cE}{\mathcal{E}}
\newcommand{\cC}{\mathcal{C}}
\newcommand{\ind}[1]{\mathbf{1}_{[#1]}}
\newcommand{\argmax}{\operatorname*{arg\,max}}
\newcommand{\argmin}{\operatorname*{arg\,min}}
\newcommand{\Tr}{\mathrm{Tr}}
\newcommand{\Situs}{\mathrm{Situs}}
\newcommand{\Cercis}{\mathrm{Cercis}}
\newcommand{\id}{\mathrm{id}}
\newcommand{\supp}{\mathrm{supp}}
\newcommand{\spec}{\mathrm{spec}}
\newcommand{\braopket}[3]{\langle#1|#2|#3\rangle}
\newcommand{\fid}[2]{F\!\left(#1,#2\right)}
\newcommand{\tdist}[2]{\frac12\|#1-#2\|_1}
\newcommand{\Holevo}{\chi}
\newcommand{\D}{\mathrm{D}}

\title{\bfseries Quantum SCX:\\
No-Cloning Strengthening of Theorem~3,\\
Entanglement Audit Advantage, and the\\
Quantum Cercis Uncertainty Relation}

\author{SCX}

\date{\today}

\begin{document}
\maketitle

\begin{abstract}
The \textsf{SCX} framework has been built entirely on classical information theory:
data are classical bits, and audit is a computation over classical statistics.
We extend \textsf{SCX} to the quantum domain by asking: if audited data are encoded as
quantum states, or if the auditor can access data through quantum channels, does
Theorem~3's noise--difficulty unidentifiability barrier still hold?  We prove
three results.  (i)~\textbf{Quantum-Strengthened Unidentifiability}: when the
auditor accesses data only through a quantum channel with bounded fidelity
$F_{\max}<1$, the indistinguishability gap $\epsilon(F_{\max})$ scales as
$O(1-F_{\max})$ --- the lower the channel fidelity, the more robust the Theorem~3
barrier.  This is \textbf{rigorously derivable} from the quantum no-cloning
theorem together with Theorem~3 and the quantum Fano inequality.
(ii)~\textbf{Entanglement Audit Advantage}: we formulate a precise conjecture
that when pre-shared entanglement $\Phi_{AB}$ is available, the auditor can
construct a quantum statistic $T_Q$ that distinguishes noise-dominated from
difficulty-dominated quantum states even when classical audit statistics
cannot.  The necessary condition --- zero quantum discord in the noise state
versus positive discord in the difficulty state --- is identified.  The conjecture
is an \textbf{open problem} requiring explicit construction of the separating
state pair.  (iii)~\textbf{Quantum Cercis Uncertainty Relation}: if the quality
and noise observables fail to commute ($[\hat{Q},\hat{N}]\neq0$), then
$\Delta Q\cdot\Delta N \geq \frac12|\langle\hat{C}\rangle|$ --- a rigorous
consequence of the Robertson--Schr\"odinger uncertainty principle.  This
constitutes an \textit{audit uncertainty principle}: greater precision in
quality estimation forces greater uncertainty in noise estimation, and vice
versa.
\end{abstract}

% ============================================================
\section{Introduction}
% ============================================================

Theorem~3 of the \textsf{SCX} framework~\citep{scx2024cercis} establishes, within
classical information theory, that noise and difficulty are unidentifiable from
Cercis Score observations alone.  The proof constructs two classical worlds that produce identical joint distributions over all \textsf{SCX} observables.  This is a \emph{statistical} indistinguishability: no classical hypothesis test has power exceeding its significance level.

Quantum mechanics introduces an entirely new kind of information barrier: the
\textbf{no-cloning theorem}~\citep{wootters1982single} asserts that an unknown
quantum state cannot be perfectly copied.  If a data holder encodes the dataset
as quantum states, the auditor \emph{cannot} make a copy to inspect --- any
measurement necessarily disturbs the system.  This suggests that quantum
mechanics might \emph{strengthen} Theorem~3 by adding a physical
unclonability layer on top of the statistical unidentifiability layer.

However, quantum mechanics also provides a countervailing resource:
\textbf{entanglement}.  Pre-shared Bell pairs enable quantum teleportation,
superdense coding, and remote state preparation --- capabilities with no
classical analogue~\citep{nielsen2010quantum}.  If entanglement allows the
auditor to extract information that classical channels cannot, then quantum
auditing might \emph{break} Theorem~3's barrier through physical means, not
merely through the cognitive lens-shift studied in HC-Theorem.

The Q-Theorem formalizes this tension.  Its central question is: \textbf{does
quantum information strengthen or weaken the SCX audit barrier?}  The answer,
developed below, is that it does \emph{both}: no-cloning strengthens the
barrier (part~i), entanglement may break it (part~ii), and the two effects
compete within a rigorous uncertainty bound (part~iii).

\section{Preliminaries}

\begin{definition}[Quantum Data Encoding]
\label{def:quantum-encoding}
A dataset $\cD=\{(x_i,y_i)\}_{i=1}^n$ is encoded into a quantum state $\rho_{\cD}$ on a finite-dimensional Hilbert space $\cH_D$. The auditor accesses data through a quantum channel $\cE: \cL(\cH_D) \to \cL(\cH_A)$.
\end{definition}

\begin{definition}[Quantum Cercis Operator]
\label{def:quantum-cercis}
The quantum Cercis operator is $\hat{S} = \hat{Q} + \eta \hat{N}$, where $\hat{Q}$ and $\hat{N}$ are observables. In general, $[\hat{Q},\hat{N}]\neq0$: quality and noise are incompatible observables.
\end{definition}

\begin{definition}[Channel Fidelity]
\label{def:channel-fidelity}
$F_{\max} = \sup_{\rho} F(\cE(\rho),\rho)$. A perfect channel has $F_{\max}=1$; any physical channel constrained by no-cloning has $F_{\max}<1$.
\end{definition}

\section{Main Results}

\subsection{Quantum-Strengthened Unidentifiability}

\begin{theorem}[Quantum-Strengthened Theorem~3]
\label{thm:quantum-strong}
For a quantum channel $\cE$ with fidelity $F_{\max}<1$, the observable Cercis expectation difference satisfies
\begin{equation}\label{eq:quantum-barrier}
    \bigl|\Tr(\hat{S}\,\cE(\rho_{\mathrm{noise}})) - \Tr(\hat{S}\,\cE(\rho_{\mathrm{hard}}))\bigr|
    \;\leq\; \epsilon(F_{\max}),
\end{equation}
where $\epsilon(F_{\max}) = (1+\eta)\sqrt{2(1-F_{\max})} + O((1-F_{\max})^{3/2})$.
\end{theorem}

\begin{proof}[Proof sketch]
The proof proceeds in three steps: (i) quantum Fano inequality, (ii) Holevo bound with channel fidelity constraint, (iii) explicit construction of the confusion strategy. The quantum Fano inequality bounds the auditor's ability to distinguish $\rho_{\mathrm{noise}}$ from $\rho_{\mathrm{hard}}$ through $\cE$. The Holevo bound constrains the accessible information by the channel fidelity $F_{\max}$, yielding the $\epsilon(F_{\max})$ expression.
\end{proof}

\subsection{Entanglement Audit Advantage Conjecture}

\begin{conjecture}[Entanglement Audit Advantage]
\label{conj:entanglement}
When pre-shared entanglement is available and the noise state has zero quantum discord while the difficulty state has positive discord, the auditor can construct a quantum statistic $T_Q$ achieving strictly super-classical discrimination power.
\end{conjecture}

\subsection{Quantum Cercis Uncertainty Relation}

\begin{theorem}[Quantum Cercis Uncertainty]
\label{thm:uncertainty}
If $[\hat{Q},\hat{N}] = i\hat{C}$, then
\[
\Delta Q \cdot \Delta N \geq \frac12 |\langle\hat{C}\rangle|,
\]
a fundamental precision trade-off between quality and noise estimation.
\end{theorem}

% ============================================================
\section{Conclusion}
% ============================================================

Q-Theorem extends \textsf{SCX} to the quantum domain: no-cloning strengthens Theorem~3's barrier, entanglement may break it, and non-commuting quality/noise observables impose a fundamental uncertainty trade-off. These results establish that quantum information both constrains and enables \textsf{SCX} auditing in ways that have no classical analogue.

% ============================================================
\bibliographystyle{plainnat}
\begin{thebibliography}{10}

\bibitem{scx2024cercis}
SCX Working Group.
\newblock ``The \textsf{SCX} Axiomatic Framework: Cercis, Situs, and Tiresias Operators,''
\newblock \emph{Technical Report}, 2024.

\bibitem{wootters1982single}
W.~K.~Wootters and W.~H.~Zurek.
\newblock ``A single quantum cannot be cloned,''
\newblock \emph{Nature}, 299:802--803, 1982.

\bibitem{nielsen2010quantum}
M.~A.~Nielsen and I.~L.~Chuang.
\newblock \emph{Quantum Computation and Quantum Information}.
\newblock Cambridge University Press, 2010.

\end{thebibliography}

\end{document}
